\documentclass[polish,polish,a4paper]{article}
\usepackage[T1]{fontenc}
\usepackage[utf8]{inputenc}
\usepackage{babel}
\usepackage{pslatex}
\usepackage{graphicx}
\usepackage{tikz}
\usepackage{pgfplots}
\usepackage{anysize}
\usepackage{pgfgantt}
\usepackage{tabularx}
\usepackage{float}
\usepackage{latexsym,amsmath}
\marginsize{2.5cm}{2.5cm}{3cm}{3cm}

\newcommand{\name}[1]{\sffamily\bfseries\scriptsize #1}
\newcommand{\frontpage}[8]{
%% #1 - nazwa kursu
%% #2 - kierunek 
%% #3 - termin 
%% #4 - temat 
%% #5 - problem
%% #6 - skład grupy
%% #7 - nr grupy
%% #8 - data

%\hfill\includegraphics[width=4cm]{PWr.png}
\vspace{2cm}

\begin{tabular}{|p{0.72\textwidth}|p{0.28\textwidth}|}
\hline
\multicolumn{2}{|c|}{}\\
\multicolumn{2}{|c|}{{\LARGE #1}}\\
\multicolumn{2}{|c|}{}\\
\hline
\name{Kierunek} & \name{Termin}\\
\multicolumn{1}{|c|}{\textit{#2}} & \multicolumn{1}{|c|}{\textit{#3}} \\
\hline
\name{Temat} & \name{data}\\
\multicolumn{1}{|c|}{\textit{#4}} & \multicolumn{1}{|c|}{\textit{#5}} \\
\hline
\name{Grupa oddająca projekt} & \name{Nr indeksu}\\
\multicolumn{1}{|c|}{\textit{#6}} & \multicolumn{1}{|c|}{\textit{#7}} \\
\hline
\end{tabular}
}

\usepackage{listings}
\usepackage{xcolor} % for setting colors

% set the default code style
\lstset{
    frame=tb, % draw a frame at the top and bottom of the code block
    tabsize=4, % tab space width
    showstringspaces=false, % don't mark spaces in strings
    numbers=left, % display line numbers on the left
    commentstyle=\color{green}, % comment color
    keywordstyle=\color{blue}, % keyword color
    stringstyle=\color{red} % string color
}

\title{Sprawozdanie SPD}
\pgfplotsset{compat=1.16}
\begin{document}
% #1 - nazwa kursu #2 - kierunek  #3 - termin #4 - temat #5 - data #6 - skład grupy #7 - nr grupy
\frontpage{Układy cyfrowe i systemy wbudowane}{Informatyka}{Poniedziałek 11:05}{Gra pamięciowa}{\today}{Paweł Sagan, Patryk Wlazłyń}{241250, 241249}
\pagestyle{}
\newpage
\tableofcontents
\newpage
\section{Wstęp}
Celem projektu było zaprojektowanie oraz wykonanie programu dedykowanego pod układ FPGA, a następnie jego przetestowanie na układzie Spartan-3E.
Wybrany przez nas temat to gra pamięciowa, w której gracz musi zapamiętać umiejscowione liczby na ekranie.
Gdy gracz naciśnie pierwszą z liczb, reszta zamienia się w niezroróżnialne pola (wypełnione prostokąty).
Celem gracza jest zapamiętanie układu liczb, oraz kliknięcie pól w kolejności jaką reprezentowały wcześniej wyświetlone liczby.

\section{Założenia projektowe}
Poniżej lista zawierająca przedstawione założenia projektowe:
\begin{itemize}
  \item Gra dedykowana jest dla jednego gracza
  \item Podstawowy interfejsem gracza jest mysz podłączona za pomocą złącza ps2
  \item Aktualny stan gry oraz komunikaty wyświetlane są na ekranie którego rozdzielczość renderowanego obrazu wynosi 48 linii po 20 znaków
  \item Początkowe umiejscowienie liczb na ekranie, losowane jest przy użyciu algorytmu pseudolosowego LCG.
  \item Pierwsza poprawną liczbą jest wartość najmniejsza spośród wyświetlonych.
  \item Ukrycie liczb następuje w momencie kliknięcia na pierwszą z nich.
  \item Odkrycie liczby większej o dwa od poprzednio odkrytej liczby, powoduje zakończenie gry oraz wyświetleniu komunikatu o przegranej.
  \item Użytkownik wygrywa gdy wszystkie liczby zostaną poprawnie odkryte. Wówczas zostaje wyświetlony stosowny komunikat.
\end{itemize}

\section{Opis projektu}
Program składa się z modułów wykonywanych w języku VHDL, które następnie połączone zostały za pomocą głównego schematu projektu.
Oprócz modułów wykonanych na potrzebę tego programu, użyte zostały komponenty dostarczone na stronie kursu, konkretnie PS2\_MOUSE oraz VGAtxt48x20.
Ze względu na ograniczony dostęp do sprzętu, dla każdego komponentu utworzony został test, który sprawdza jego poprawną implementacje.
Dodatkowo aby zminimalizować ilość błędów, komponenty zostały pogrupowane w zbiory zawierające najmniejsza możliwą ilość modułów stanowiącą logiczną całość.
Dzięki takiemu podejściu analiza błędów syntezy jest mniej złożona.
Projekt dedykowany jest pod zestaw Digilent S3E-Starter, który to posiada wyprowadzone złącze PS2 oraz konektor VGA.
\newline
Głównym komponentem programu jest moduł logic który spaja ze sobą wszystkie stworzone moduły.
W poniższej sekcji znajduje się dokładny opis zaprojektowanych modułów oraz relacji pomiędzy nimi. 
\newpage
\subsection{Wykorzystane moduły}
\subsubsection{Stworzone moduły}
\begin{itemize}
 \item
  Licznik - Prosty 16-bitowy moduł zliczający cykle zegara oraz podający informacje o aktualnym stanie licznika na magistrale.
  Użyty jako "seed" dla modułu implementującego generowanie liczb pseudolosowych.

 \item
  PS2\_Controler - Moduł służy do wyliczania aktualnej pozycji myszki.
  Wyliczone dane ograniczone są przez minimalną oraz maksymalną rozdzielczość ekranu.
  Oznacza to, że moduł nigdy nie zwróci ujemnej pozycji oraz pozycji większej niż założona.
  Układ propaguje również informacje o kliknięciu przez użytkownika lewego przycisku myszy.
  Informacja widnieje na linii tak długo, aż inny moduł nie poprosi o jej wyczyszczenie.
  Wystawione dane pozycji myszki w momencie kliknięcia, pozostają niezmienione, aż do momentu ustawienia linii na stan niski. 

 \item
  vga\_driver - Prosty układ potrafiący przekazywać dane od logiki programu na wejście sterownika VGAtxt48x20.
  Układ posiada 4 tryby.
  Pierwszy tryb umożliwia renderowanie pamięci bez ukrywania jej zawartości.
  Oznacza to, że wszystkie wygenerowane liczby renderowane są jako ich odpowiadająca znakowa reprezentacja.
  Kolejny tryb umożliwia renderowanie ukryte.
  Oznacza to, że wszystkie liczby w pamięci wyświetlane są jako białe kafelki.
  Dwa ostanie tryby umożliwiają wyświetlenie informacji o stanie gry użytkownika.
  Pierwszy z nich renderuje informacje o przegranej użytkownika, drugi zaś o wygranej.
  Moduł przyjmuje informacje o pozycji myszki, dzięki czemu jest w stanie wygenerować symbol myszy w odpowiednim miejscu.

 \item
  lcg - Moduł odpowiedzialny za generowanie liczb pseudolosowych.
  Wewnętrznie używa implementacji algorytmu LCG (Linear congruential generator).
  Parametry algorytmu zostały dobrane eksperymentalnie, bazując na rozwiązaniach w bibliotece glibc (standardowa biblioteka C, GNU).
  Dodatkowo dobrane parametry zostały zasymulowane w języku Python, aby potwierdzić równomierny rozkład losowanych liczb.

 \item
  object\_generator - Moduł odpowiedzialny za wygenerowanie kolejnych obiektów wyświetlanych na ekranie w losowych miejscach.
  Wewnetrzna implementacja korzysta z zmodyfikowanego algorytmu "Fisher–Yates shuffle".
  Aby powierdzić równomierny rozkład podczas operacji mieszania, zmodyfikowany algorytm został w pierwszej kolejności zaimplementowany w języku Python.
  Modyfikacja polega na poszerzeniu zakresu losowania elementów, które zamieniane są miejscami.
  Dzięki temu dzielenie przy obliczaniu reszty z dzielenia, wykorzystuje dzielenie przez stałą, przez to nie jest wymagany układ dzielący.
  object\_generator operuje na pamięci ram w celu wygenerowania losowego położenia obiektów.
  Parametr wielkości miejsc pustych nie może być zmieniony dynamicznie (w trakcie działania układu), natomiast jest łatwo modyfikowalny w kodzie.
  Wymaga to jednak ponownej syntezy.
  Parametr ilości generowanych obeiktów zajmujących nie puste pola jest natomiast całkowicie konfigurowalny w czasie działania układu.

 \item
   ram16 - prosty moduł implementujący pamięć RAM z 16-bitową szyną adresową dla 8-bitowych komórek danych.

 \item
   memmux - prosty moduł implementujący multiplexer dostępu do modułu ram16.

 \item
  logic - główny moduł spajający i zarządzający.
  Implementuje główną logikę gry, oraz zarządzanie jej stanem.
  Steruje restartowaniem oraz inicjalizacją gry, przełącza moduł vga\_driver w odpowiedni tryb wyświetlania, obsługuje zdarzenia generowane przez moduł 
  PS2\_Controler, oraz zarządznie dostępem do współdzielonej magistarali do pamięci RAM za pomocą multiplexera memmux.


\end{itemize}

\subsubsection{Zewnętrzne moduły}
\begin{itemize}
 \item
  Moduł PS2\_Mouse - Moduł ten przekazuje dane dostarczone przez aktualnie podłączoną myszkę.
  Kiedy dane są aktywne układ ustawia jedno cyklowy stan wysoki na linii.
  Oprócz tego dostarcza informacje o przesunięciu myszki względem ostatniej pozycji w osi X oraz Y.
  Oprócz tego moduł przekazuje informacje na temat zwrotu myszki w poszczególnych kierunkach oraz informacje na temat wciśniętych przycisków.
  Oś X, Y oraz status są 8 bitowymi magistralami.

 \item
  Moduł VGAtxt48x20 - moduł wyświetla przekazane dane na ekranie monitora.
  Dane podawane są na 8 bitową magistrale.
  Moduł wyświetla znaki zakodowane według standardowej tablic ASCII, wyjątkiem są wartości z przedziału 0x0 - 0x1F.
  Informacje na temat wspominanego zakresu można znaleźć w dokumentacji modułu.
  Po wyświetleniu zadanego znaku, moduł automatycznie zwiększa pozycję na której wystąpi wyświetlenie kolejnego znaku.
  Możliwy jest również powrót kursora do początku linii lub początku ekranu.

\end{itemize}

\subsection{Osiągnięte cele}
Wszystkie moduły do zrealizowania kompletnej logiki gry zostały utworzone oraz przetestowane przy użyciu symulatora.
Proces generowania bit-streamu oraz syntezy przebiegł pomyślnie.
Ze względu na ograniczony dostęp do sprzętu niemożliwe jest przetestowanie układu na sprzęcie.

\section{Instrukcja obsługi}
Po załadowaniu programu na urządzenie, gracz rozpoczyna nową grę przez kliknięcie w dowlnym miejscu na ekranie.
Następnie stan nowej gry zostaje zainicjalizowany w pamięci, oraz zostają wylosowane miejsca na kolejne liczby.
Liczby zostają ukazane na ekranie, a zadaniem gracza jest kliknięci ich w kolejności reprezentowanej przez wartości liczb.
Kliknięcie chociaż jedej z nich w złej kolejności spowoduje przegraną (aby aby rozpocząć na nowo, należy kliknąć w dowolne miejsce na ekranie).
Aby wygrać wszystkie obiekty muszą zostać odkryte.

\section{Podsumowanie}
Projektowanie aplikacji oraz tworzenie ich w języku VHDL wymaga podejścia innego niż w przypadku programownaia w językach wysokiego poziomu, które
wykonywane są przez procesory ogólnego przeznaczenia.
Należy pamiętać o zależnościach czasowych oraz o tym, że łączenie instrukcji oferowanych przez język VHDL nie odzwierciedla zachowania znanych z innych wysokopoziomowych języków programowania, jak np. to, że rezultaty niektórych przypisań widoczne są dopiero z następnym cyklem zegara.

\end{document}
